\documentclass[12pt]{article}

\usepackage{sbc-template}
\usepackage{graphicx,url}
\usepackage[utf8]{inputenc}
\usepackage[english]{babel}
\usepackage{amsmath}
\usepackage{booktabs}
\usepackage{subcaption}
\usepackage{float}
     
\sloppy

\title{Generative Artificial Intelligence and Software Engineering: Challenges and Opportunities.}

\author{Gabriel Ferreira Silva\inst{1}}

\address{Instituto de Informática -- Universidade Federal de Goiás (UFG)\\
  Goiânia -- GO -- Brasil
  \email{ferreira.silva@discente.ufg.br}
}

\begin{document} 

\maketitle

\begin{abstract}
Requirements elicitation remains time-consuming and susceptible to information loss. This work proposes a multi-agent system powered by generative AI to automate the pipeline from meeting audio to structured user stories. The system comprises three specialized agents: transcription (Whisper), persona identification and user story generation (GPT-4o-mini), and requirements formatting. Evaluation using the AMI Meeting Corpus demonstrated 100\% persona identification accuracy (28/28 across 7 meetings) and 83\% INVEST compliance. The qualitative framework yielded 9/10, validating GenAI-powered requirements automation feasibility. This research provides modular architecture, empirical validation, and open-source implementation for future deployment.
\end{abstract}


\section{Introduction}
\label{sec:introduction}

Requirements engineering constitutes a fundamental phase in software development, serving as the bridge between stakeholder needs and technical implementation. Traditional requirements elicitation processes, however, remain predominantly manual, time-consuming, and susceptible to information loss and inconsistencies. Organizations typically conduct extensive stakeholder meetings to gather requirements, followed by manual transcription, analysis, and documentation—a workflow that demands substantial time investment from requirements analysts and often results in incomplete or ambiguous requirements specifications.

The emergence of Large Language Models (LLMs) and generative artificial intelligence presents transformative opportunities for automating and enhancing requirements engineering workflows. Recent advances in natural language processing, particularly through models such as GPT-4 and Whisper, demonstrate unprecedented capabilities in understanding, analyzing, and generating human language across diverse contexts. These capabilities suggest potential for automating critical requirements engineering tasks, including meeting transcription, stakeholder analysis, and requirements documentation.

Despite growing interest in applying LLMs to requirements engineering, existing research predominantly addresses isolated subtasks rather than complete end-to-end workflows. Current work focuses separately on automated transcription, requirements classification, or user story generation, but limited research explores integrated systems that transform raw meeting audio into structured, stakeholder-specific requirements documentation. Furthermore, while multi-agent architectures demonstrate advantages in other software engineering domains through modularity and task specialization, their systematic application to requirements engineering remains underexplored.

This work proposes, implements, and evaluates a multi-agent system that automates this complete pipeline, from meeting audio recordings to structured user stories with explicit persona identification. The research employs the AMI Meeting Corpus for validation and adopts a qualitative assessment framework to evaluate persona identification accuracy and user story quality.

\subsection{Research Objectives}

The primary objective of this research is to investigate the feasibility of using generative AI-powered multi-agent systems for automated requirements elicitation from conversational audio. Specifically, this work aims to:

\begin{enumerate}
    \item Design and implement a modular multi-agent architecture for automated requirements engineering.
    
    \item Develop an audio transcription agent capable of processing meeting recordings in multiple languages.
    
    \item Create an identification agent that extracts stakeholder personas and generates structured user stories.
    
    \item Establish a qualitative evaluation framework for assessing automatically generated requirements artifacts.
    
    \item Identify limitations and opportunities for future quantitative validation.
\end{enumerate}

\subsection{Main Contributions}

This research provides the following contributions:

\begin{enumerate}
    \item \textbf{Novel multi-agent architecture}: A modular system design with three specialized agents (transcription, identification, formatting) for automated requirements elicitation.
    
    \item \textbf{End-to-end automation prototype}: A functional implementation demonstrating the complete pipeline from meeting audio to structured user stories.
    
    \item \textbf{Persona-centric requirements extraction}: An automated approach to identifying stakeholders and generating persona-specific user stories.
    
    \item \textbf{Qualitative evaluation methodology}: A systematic assessment framework achieving 9/10 overall quality score and 100\% persona identification accuracy (28/28 across 7 meetings).
    
    \item \textbf{Open-source implementation}: Complete source code available on GitHub (\url{https://github.com/gabrielsfg/PFC1}).
\end{enumerate}

\subsection{Document Organization}

The remainder of this paper is organized as follows. Section~\ref{sec:related} reviews relevant literature on large language models in requirements engineering, automated user story generation, multi-agent systems, speech recognition technologies, and evaluation methodologies, identifying key research gaps that motivate this work. Section~\ref{sec:method} presents the research methodology, including problem definition, dataset selection, multi-agent architecture design, and detailed implementation of each agent. Section~\ref{sec:results} reports experimental results and qualitative evaluation of the implemented agents, including persona identification accuracy, user story quality assessment, and INVEST criteria compliance analysis. Section~\ref{sec:conclusion} summarizes the main contributions, acknowledges current limitations, and outlines directions for future work in the second phase of this project (PFC2), including completion of the Formatting Agent and comprehensive quantitative validation.


\section{Related Work}
\label{sec:related}

The integration of Large Language Models (LLMs) and artificial intelligence techniques in Requirements Engineering represents a rapidly evolving research area. This section reviews relevant work across multiple dimensions: fundamental surveys on LLMs, their application in requirements engineering tasks, automated user story generation, multi-agent systems, speech recognition technologies, datasets, and prompting techniques.

\subsection{Large Language Models: Foundations and Surveys}

The emergence of LLMs has fundamentally transformed natural language processing capabilities across multiple domains. Hort et al. \cite{hort:hal-05341748} conducted a comprehensive tertiary study analyzing 274 surveys on LLMs, systematically categorizing research across model architectures, training methodologies, evaluation approaches, and application domains. Their work reveals the rapid growth of LLM research, with significant focus on foundational capabilities such as text comprehension, generation, world knowledge representation, and reasoning abilities. This survey of surveys provides essential context for understanding how LLMs have evolved from basic language modeling to sophisticated systems capable of complex reasoning and task execution, establishing the theoretical foundation for their application in specialized domains such as requirements engineering.

\subsection{Datasets for LLM-based Requirements Engineering}

The availability and quality of domain-specific datasets critically influence LLM performance in requirements engineering tasks. Motger et al. \cite{motger2025share} conducted a systematic mapping study identifying and characterizing 62 publicly available datasets used in LLM-based RE research. Their analysis revealed significant research gaps, including limited coverage for requirements elicitation tasks (only 2 datasets), scarce resources for management activities beyond traceability, and overwhelming linguistic bias toward English (58 out of 62 datasets). The study also identified dataset scale limitations, with most collections containing fewer than 1,000 artifacts, constraining their applicability for large-scale model training and fine-tuning. These findings highlight the data scarcity challenge that motivates the selection of well-established datasets such as the AMI Meeting Corpus \cite{ami_dataset} for requirements elicitation research.

\subsection{Automated User Story Generation}

Automated generation of user stories represents a critical application area for LLMs in agile requirements engineering. Delgado et al. \cite{delgado2024aidriven} investigated automated user story generation using two distinct approaches: N-gram representation with linguistic heuristics and GPT-3-based generation. Their evaluation across diverse user story corpora using BLEU, ROUGE-N, and BERTScore metrics demonstrated that GPT-3 achieved superior performance (ROUGE-N=0.46, BLEU=0.27, BERTScore=0.69) compared to traditional N-gram approaches, though both methods revealed opportunities for quality improvement. This work established baseline performance metrics for LLM-based user story generation and validated the viability of using generative models for this task.

The transformation of meeting conversations into structured user stories, however, remains underexplored. While existing work focuses primarily on generating user stories from existing textual requirements or specifications, limited research addresses the complete pipeline from audio meeting recordings through transcription, persona identification, to user story generation—the gap this work aims to address.

\subsection{User Story Quality and Structure}

The quality of user stories, whether manually or automatically generated, significantly impacts their utility in software development. Traditional quality frameworks emphasize adherence to the INVEST criteria (Independent, Negotiable, Valuable, Estimable, Small, Testable) to ensure user stories effectively guide implementation activities. Research on automated evaluation of user story quality demonstrates that LLM-generated stories often achieve high format compliance but may exhibit variability in granularity and independence, necessitating systematic quality assessment frameworks.

\subsection{Class Diagrams and Conceptual Models from User Stories}

Beyond user story generation, researchers have explored downstream requirements engineering tasks that leverage user stories as input. Ren et al. \cite{ren2024combining} investigated Chain-of-Thought (CoT) prompting techniques to automate the extraction of conceptual models from user stories. Their findings demonstrated that well-crafted few-shot prompts enabled LLMs to outperform guided human extraction in identifying classes, though qualitative analysis revealed areas of suboptimal performance in relationship identification and attribute assignment. This work illustrates both the potential and current limitations of LLMs in transforming natural language requirements into structured engineering artifacts.

\subsection{Multi-Agent Systems for Software Engineering}

While multi-agent architectures offer modularity, task specialization, and maintainability advantages, their application in automated requirements engineering workflows remains limited. Sami et al. \cite{sami2024experimentingmultiagent} explored multi-agent approaches to software development, proposing a unified platform for coordinating specialized agents in various development tasks. Their work demonstrates the potential of decomposing complex software engineering processes into collaborative agent-based workflows, though specific application to requirements elicitation and documentation remains underexplored. The modular decomposition of requirements engineering processes into specialized agents—each responsible for distinct subtasks such as transcription, analysis, and formatting—represents an architectural pattern with potential for improving system robustness, maintainability, and extensibility.

\subsection{Speech Recognition for Meeting Analysis}

Automated transcription of meeting audio represents the initial stage in audio-based requirements elicitation pipelines. Modern speech recognition systems, particularly OpenAI's Whisper model \cite{whisper2022}, demonstrate robust multilingual transcription capabilities with integrated language detection and punctuation generation. However, research specifically addressing the application of speech recognition technologies to requirements engineering meetings, including challenges related to domain-specific terminology, multiple speakers, and informal conversational contexts, remains limited in the requirements engineering literature.

\subsection{Metrics and Evaluation in Requirements Engineering}

The evaluation of automated requirements engineering tools necessitates appropriate metrics that capture both syntactic correctness and semantic quality. Traditional software engineering metrics \cite{bokhari2011metrics} provide frameworks for assessing requirements quality, but their application to LLM-generated outputs requires adaptation. The challenge of evaluating persona identification accuracy, user story format compliance, and semantic coherence in automatically extracted requirements motivates the development of specialized qualitative assessment frameworks tailored to generative AI outputs.

\subsection{Identified Gaps and Research Opportunities}

The literature review reveals several critical gaps that motivate this research:

\begin{enumerate}
    \item \textbf{End-to-End Audio-to-Requirements Pipelines}: Limited research addresses complete workflows from meeting audio recordings to structured requirements documentation, with most work focusing on individual subtasks in isolation.
    
    \item \textbf{Multi-Agent Architectures for RE}: While multi-agent systems demonstrate advantages in other software engineering domains \cite{sami2024experimentingmultiagent}, their systematic application to requirements engineering workflows remains underexplored.
    
    \item \textbf{Persona-Centric Requirements Extraction}: Automated identification of distinct stakeholder personas from conversational text and their linkage to specific requirements represents a novel contribution not adequately addressed in existing literature.
    
    \item \textbf{Qualitative Evaluation Frameworks}: Systematic methodologies for qualitatively assessing LLM-generated requirements artifacts, particularly in early development stages before comprehensive quantitative validation, require further development.
    
    \item \textbf{Dataset Scarcity for Requirements Elicitation}: The identified shortage of publicly available datasets \cite{motger2025share} specifically designed for requirements elicitation tasks limits reproducibility and comparative evaluation of proposed approaches.
\end{enumerate}

This work addresses these gaps by proposing, implementing, and qualitatively evaluating a multi-agent system that automates the complete pipeline from meeting audio to structured user stories, with explicit persona identification and systematic quality assessment.

\section{Research Method}
\label{sec:method}

This section presents the methodology adopted to develop and evaluate the proposed multi-agent system. The following subsections detail: (i) the problem definition, (ii) the dataset selected for validation, (iii) the system architecture, and (iv) the implementation of each agent.


\subsection{Problem Definition}
\label{sec:problem}

Traditional requirements elicitation relies on manual processes involving:

\begin{enumerate}
    \item Conducting meetings with stakeholders to gather requirements
    \item Manual transcription of meeting recordings or note-taking during discussions
    \item Analysis of transcripts to identify different stakeholders (personas)
    \item Extraction and documentation of requirements as user stories
    \item Manual formatting and organization of user stories following established standards
\end{enumerate}

This manual approach presents several critical challenges:

\begin{itemize}
    \item \textbf{Time consumption}: Transcription and analysis are labor-intensive tasks that significantly extend project timelines
    \item \textbf{Information loss}: Important details and context may be missed during manual note-taking or transcription
    \item \textbf{Inconsistency}: Manual documentation often results in inconsistent formatting and varying levels of detail across user stories
    \item \textbf{Cognitive overload}: Requirements engineers must simultaneously participate in discussions while capturing information accurately
    \item \textbf{Delayed feedback}: The gap between meeting occurrence and documented requirements can lead to stakeholder memory decay
\end{itemize}

Recent advances in Generative Artificial Intelligence, particularly Large Language Models (LLMs), present an opportunity to automate significant portions of this workflow. However, applying GenAI to requirements engineering faces specific challenges:

\begin{itemize}
    \item \textbf{Context understanding}: Accurately identifying different stakeholders and their specific needs from conversational text
    \item \textbf{Structured output generation}: Converting natural conversation into properly formatted, actionable user stories
    \item \textbf{Quality assurance}: Ensuring generated requirements meet software engineering standards (e.g., INVEST criteria)
\end{itemize}

This work addresses these challenges by proposing an automated multi-agent system that transforms audio recordings from requirements gathering meetings into structured, well-formatted user stories, reducing manual effort while maintaining quality and completeness.


\subsection{Dataset}
\label{sec:dataset}

To validate and test the proposed system, we utilize the \textbf{AMI Meeting Corpus} \cite{ami_dataset}, a publicly available dataset developed by the European AMI consortium (FP6-506811). The dataset is accessible through both the HuggingFace platform and the official University of Edinburgh repository.

\subsubsection{Dataset Characteristics}

The AMI Meeting Corpus consists of approximately 100 hours of meeting recordings, totaling 279 dialogues divided into training, validation, and test sets. The dataset presents the following technical characteristics:

\begin{itemize}
    \item \textbf{Volume}: Approximately 100 hours of recorded meetings
    \item \textbf{Distribution}: 279 dialogues (209 training, 42 validation, 28 test)
    \item \textbf{Language}: Predominantly English, with majority of non-native speakers
    \item \textbf{Speaker identification}: Each participant has unique speaker labels
    \item \textbf{Context}: Professional meeting environment conversations
    \item \textbf{Transcriptions}: High-quality manual orthographic transcriptions with word-level timestamps
    \item \textbf{Annotations}: Dialogues annotated with dialogue acts, topic segmentation, summaries, and named entities
\end{itemize}

\subsubsection{Corpus Composition}

The dataset is structured into two main meeting types:

\begin{enumerate}
    \item \textbf{Scenario-based meetings} (approximately 2/3 of corpus): 138 meetings where participants assume different roles in a fictional product design team, conducting a project from kickoff to completion over a workday. Roles include: industrial designer, interface designer, marketing, and project manager.
    
    \item \textbf{Natural meetings} (approximately 1/3 of corpus): 59 real meetings from various domains that would have occurred regardless of recording, providing authentic group interaction data.
\end{enumerate}

Table~\ref{tab:ami_distribution} presents the distribution of meetings across different subsets.

\begin{table}[htb]
\centering
\caption{AMI Meeting Corpus distribution}
\label{tab:ami_distribution}
\begin{tabular}{lcc}
\toprule
\textbf{Subset} & \textbf{Number of Meetings} & \textbf{Percentage} \\
\midrule
Training & 209 & 74.9\% \\
Validation & 42 & 15.1\% \\
Test & 28 & 10.0\% \\
\midrule
\textbf{Total} & \textbf{279} & \textbf{100\%} \\
\bottomrule
\end{tabular}
\end{table}

\subsubsection{Rationale for Dataset Selection}

The AMI Meeting Corpus was selected for this work due to several key advantages:

\begin{itemize}
    \item \textbf{Available ground truth}: Pre-existing speaker identification enables objective validation of persona identification performed by the system
    \item \textbf{Natural conversations}: Contains authentic dialogues in professional meeting contexts, similar to the proposed system's target application scenario
    \item \textbf{Scientific validation}: Widely used dataset in academic literature for evaluating conversation processing and meeting summarization systems
    \item \textbf{Public availability}: Licensed under Creative Commons Attribution 4.0 International (CC BY 4.0), ensuring study reproducibility
    \item \textbf{Rich annotations}: Multiple annotation layers (dialogue acts, topic segmentation, summaries) enable future validation of additional system functionalities
    \item \textbf{Structured format}: Data provided in JSON format with well-defined fields (id, dialogue, summary), facilitating automated processing
\end{itemize}

\subsubsection{Dataset Usage in the Project}

For this work, AMI Corpus dialogues are used as input to the multi-agent system. The utilization process involved:

\begin{enumerate}
    \item \textbf{Access}: Dataset download through HuggingFace platform using the \texttt{datasets} library
    \item \textbf{Preprocessing}: Extraction of \texttt{dialogue} and \texttt{id} fields from each example
    \item \textbf{Validation}: Comparison of persona identification performed by the Identification Agent with original corpus speaker labels
\end{enumerate}

The choice to use existing transcriptions rather than raw audio allowed initial validation to focus on the system's ability to identify personas and extract user stories from text, isolating this stage from potential errors introduced by the transcription process.


\subsection{Multi-Agent System Architecture}
\label{sec:architecture}

The proposed system operates through a sequential processing pipeline composed of three specialized agents.

\subsubsection{System Workflow}

The system operates following a sequential processing flow:

\begin{enumerate}
    \item The \textbf{Transcription Agent} captures audio and generates text files (.txt) with transcribed content
    \item The \textbf{Identification Agent} processes transcriptions, identifies personas, and extracts user stories
    \item The \textbf{Formatting Agent} standardizes user stories into structured format
\end{enumerate}

\subsubsection{Inter-Agent Communication}

Agents communicate through a shared file system, where each agent monitors the output of the previous agent, ensuring automated and continuous processing. This design provides:

\begin{itemize}
    \item \textbf{Loose coupling}: Agents operate independently, enabling parallel development and testing
    \item \textbf{Fault tolerance}: Failure in one agent does not necessarily cascade to others
    \item \textbf{Scalability}: New agents can be added to the pipeline without modifying existing ones
    \item \textbf{Monitoring capability}: File-based communication enables easy tracking and debugging of pipeline execution
\end{itemize}

The file monitoring mechanism is implemented using the \textbf{watchdog} library, which continuously monitors output directories. When a new file is detected, the subsequent agent is automatically triggered, guaranteeing continuous and automated processing flow in the multi-agent system.


\subsection{Agent 1: Audio Transcription}
\label{sec:agent1}

The first agent in the system is responsible for capturing audio conversations and automatically converting them to text, providing transcriptions that will be analyzed by the Persona Identification Agent.

\subsubsection{Functionality}

The agent offers two interaction interfaces: a graphical user interface (GUI) and a command-line interface (CLI), allowing usage flexibility according to user preference. When the user initiates a recording, the system captures microphone audio and, upon completion, automatically saves the file to the \texttt{data/input/} directory.

The transcription process is automatically triggered through a file monitoring system. As soon as new audio is detected, the agent immediately initiates transcription using OpenAI's Whisper model \cite{whisper2022}, without requiring manual intervention.

\subsubsection{Transcription Technology}

For audio-to-text conversion, the agent utilizes the \textbf{faster-whisper} library, an optimized implementation of OpenAI's Whisper model based on the CTranslate2 framework \cite{ctranslate2}. This implementation was chosen for offering superior processing speed and lower memory consumption compared to original Whisper, while maintaining transcription quality.

The system uses by default the Whisper \texttt{small} model, which offers a good balance between speed and transcription quality.

\subsubsection{Output Structure}

The agent generates two file types in the \texttt{data/output/} directory:

\begin{enumerate}
    \item \textbf{TXT file}: Contains complete transcription in continuous text
    \item \textbf{JSON file}: Contains structured data including detected language, detection probability, audio duration, and individual segments with timestamps
\end{enumerate}

These files serve as input for the Persona Identification Agent, which will process the content to extract personas and user stories.

\subsubsection{Pipeline Automation}

The file monitoring mechanism described in Section~\ref{sec:architecture} automatically triggers the Identification Agent when new transcription files are detected.


\subsection{Agent 2: Persona Identification and User Story Extraction}
\label{sec:agent2}

The second agent is responsible for processing transcriptions generated by the previous agent, automatically identifying personas present in conversations and extracting relevant user stories.

\subsubsection{Agent Objectives}

This agent's main objectives are:

\begin{itemize}
    \item Identify personas present in the transcribed conversation
    \item Extract relevant information from each persona: name/identifier, role/function, characteristics, and context
    \item Generate user stories following INVEST criteria
    \item Generate user stories in the format: "As a [persona], I want [goal] so that [benefit/reason]"
\end{itemize}

\subsubsection{Technologies and Libraries}

The agent was implemented in Python 3.11.9, utilizing:

\begin{itemize}
    \item \textbf{openai (v1.54.3)}: Interface for GPT-4o-mini model access
    \item \textbf{python-dotenv}: Secure API key management
    \item \textbf{pydantic}: JSON output validation and structuring
    \item \textbf{tenacity}: Retry strategies for API resilience
\end{itemize}

\subsubsection{AI Model Selection}

The \textbf{GPT-4o-mini} model from OpenAI was chosen for the following reasons:

\begin{itemize}
    \item Adequate cost-benefit ratio for processing large text volumes
    \item Contextual understanding capability of conversations
    \item Native support for generating structured JSON outputs
    \item Processing speed suitable for real-time application
\end{itemize}

\subsubsection{Identification Process}

The agent operates through an automated pipeline:

\begin{enumerate}
    \item \textbf{Detection}: New .txt files trigger automatic processing
    \item \textbf{Prompt Engineering}: Text is structured into a prompt instructing the model to identify personas and generate user stories following INVEST criteria
    \item \textbf{Processing}: GPT-4o-mini analyzes the transcription
    \item \textbf{Validation}: Pydantic validates the JSON structure
    \item \textbf{Persistence}: Results are saved to the output directory
    \item \textbf{Retry}: Up to 3 attempts on API failures
\end{enumerate}

\subsubsection{Output Format}

The agent generates structured JSON files containing identified personas (with name, role, characteristics, and context), associated user stories, and processing metadata.


\section{Results and Discussion}
\label{sec:results}

\subsection{Agent 1: Audio Transcription Results}
\label{sec:results_agent1}

The Transcription Agent was evaluated through preliminary tests using 5 audio recordings with durations ranging from 5 seconds to 1 minute, in both English and Portuguese. Table~\ref{tab:agent1_tests} presents the test configuration.

\begin{table}[H]
\centering
\caption{Agent 1 preliminary test configuration}
\label{tab:agent1_tests}
\begin{tabular}{lcc}
\toprule
\textbf{Test Parameter} & \textbf{Range/Values} & \textbf{Results} \\
\midrule
Number of audio samples & 5 & All processed \\
Duration range & 5s -- 60s & All durations \\
Languages tested & English, Portuguese & Both detected \\
Audio quality levels & Low, Medium, High & All transcribed \\
\bottomrule
\end{tabular}
\end{table}

\subsubsection{Language Detection Accuracy}

The agent successfully identified the correct language in all 5 test cases (100\% accuracy), demonstrating reliable language detection capability for both English and Portuguese audio inputs. This performance validates the effectiveness of the Whisper model's built-in language identification feature.

\subsubsection{Transcription Quality}

Qualitative analysis of the generated transcriptions revealed consistent accuracy across different audio quality levels. The agent produced transcriptions with appropriate punctuation and, for Portuguese texts, correct accentuation (e.g., ``é'', ``ã'', ``ç''). Transcription quality remained stable regardless of audio duration or recording quality, suggesting robust performance of the faster-whisper implementation.

\subsubsection{Current Limitations}

As this agent is still in early development stages, comprehensive quantitative evaluation was not performed. The preliminary tests focused on functional validation rather than detailed performance metrics. Future work (PFC2) will include:

\begin{itemize}
    \item Word Error Rate (WER) calculation
    \item Comparative analysis with ground truth transcriptions
    \item Processing time benchmarks
    \item Evaluation with larger and more diverse audio samples
    \item Testing with multiple speakers and background noise conditions
\end{itemize}

The current qualitative assessment indicates that the agent meets basic functional requirements for audio transcription, providing a suitable foundation for the subsequent persona identification stage.

\subsection{Agent 2: Persona Identification and User Story Generation Results}
\label{sec:results_agent2}

The Identification Agent was evaluated using 7 dialogues randomly selected from the AMI Meeting Corpus. All selected dialogues were scenario-based meetings where participants assumed specific roles in a fictional product design team. Table~\ref{tab:agent2_results} presents the quantitative overview of the generated outputs.

\begin{table}[htb]
\centering
\caption{Agent 2 output summary across 7 processed dialogues}
\label{tab:agent2_results}
\begin{tabular}{lcccc}
\toprule
\textbf{Meeting ID} & \textbf{Personas} & \textbf{User Stories} & \textbf{Format} & \textbf{Quality} \\
\midrule
meeting\_0   & 4 & 8 & Yes & Excellent \\
meeting\_2   & 4 & 4 & Yes & Excellent \\
meeting\_30  & 4 & 4 & Yes & Excellent \\
meeting\_58  & 4 & 4 & Yes & Excellent \\
meeting\_113 & 4 & 4 & Yes & Excellent \\
meeting\_117 & 4 & 4 & Yes & Excellent \\
meeting\_127 & 4 & 4 & Yes & Excellent \\
\midrule
\textbf{Average} & \textbf{4.0} & \textbf{4.6} & \textbf{100\%} & \textbf{---} \\
\bottomrule
\end{tabular}
\end{table}

\subsubsection{Persona Identification Quality}

The agent successfully identified personas in all 7 processed dialogues. Table~\ref{tab:persona_accuracy} presents a qualitative comparison between identified personas and the expected roles in AMI scenario-based meetings.

\begin{table}[htb]
\centering
\caption{Persona identification accuracy compared to AMI Corpus ground truth}
\label{tab:persona_accuracy}
\begin{tabular}{lp{4cm}p{5cm}}
\toprule
\textbf{Meeting} & \textbf{Expected Roles (AMI)} & \textbf{Identified Personas} \\
\midrule
meeting\_0 & PM, ME, UI, ID & Moderator, Marketing Expert, Designer (UI), Designer (Func) \\
\midrule
meeting\_2 & PM, ME, UI, ID & Project Coordinator, Marketing, UI Designer, Industrial Designer \\
\midrule
meeting\_117 & PM, ME, UI, ID & Meeting Moderator, Presenter, Technical Expert, Creative Collaborator \\
\bottomrule
\end{tabular}
\end{table}

Note: PM = Project Manager, ME = Marketing Expert, UI = User Interface Designer, ID = Industrial Designer

The agent demonstrated high accuracy in persona identification, successfully detecting 4 distinct personas in all 7 evaluated dialogues. Manual analysis of the meeting transcripts confirmed that each dialogue contained exactly 4 speakers, and the agent correctly identified each speaker as a unique persona, achieving 100\% detection accuracy (28/28 personas across all meetings). Furthermore, qualitative evaluation revealed that the agent consistently assigned appropriate roles to each identified persona, with role assignments aligning with the participants' actual contributions and responsibilities within the meeting context.

\subsubsection{Persona Characterization Quality}

For each identified persona, the agent consistently extracted two characteristics to support the identification of relevant user stories. Analysis of the results revealed that the agent leveraged speech patterns in each participant's contributions to infer their primary characteristics, thereby facilitating the generation of user stories aligned with their preferences and priorities. A clear example is Speaker B in meeting\_0, who repeatedly discussed the remote control's model, color, and style. Based on these recurring themes in the participant's speech, the agent correctly identified ``aesthetic-focused'' as a defining characteristic, as illustrated in Table~\ref{tab:persona_example}.

\begin{table}[htb]
\centering
\caption{Example persona characterization from meeting\_0}
\label{tab:persona_example}
\begin{tabular}{p{3cm}p{8cm}}
\toprule
\textbf{Attribute} & \textbf{Content} \\
\midrule
\textbf{ID} & persona\_2 \\
\textbf{Name} & Speaker B \\
\textbf{Role} & Designer \\
\textbf{Characteristics} & Creative, Aesthetic-focused \\
\textbf{Context} & Describes the remote control design and appearance, emphasizing form and colors \\
\textbf{Relevant Excerpts} & ``this is a model of the remote control which we are going to build.''; ``it is attractive um and it's it's blue in colour.'' \\
\bottomrule
\end{tabular}
\end{table}

The extracted characteristics demonstrated strong coherence with the actual dialogue content. The agent successfully identified personality traits and work focus areas that aligned with each participant's contributions during the meeting.

\subsubsection{User Story Generation Quality}

Generated user stories were evaluated against the INVEST criteria (Independent, Negotiable, Valuable, Estimable, Small, Testable). All generated stories followed the standard format: ``As a [persona], I want [action] so that [benefit/reason].''

Example user story from meeting\_0 (Speaker A - Designer):

\begin{quote}
\textit{``As a Designer, I want to include a mute button on the remote control so that users can quickly silence the device.''}
\end{quote}

This user story achieves a perfect INVEST score (6/6): it can be implemented independently, details are negotiable, provides clear user value, has estimable scope, is appropriately small, and is testable.

Table~\ref{tab:invest_scores} presents INVEST scores for a sample of generated user stories across different meetings.

\begin{table}[htb]
\centering
\caption{INVEST criteria evaluation for selected user stories}
\label{tab:invest_scores}
\begin{tabular}{lccccccc}
\toprule
\textbf{Story} & \textbf{I} & \textbf{N} & \textbf{V} & \textbf{E} & \textbf{S} & \textbf{T} & \textbf{Score} \\
\midrule
meeting\_0\_story\_1 & Y & Y & Y & Y & Y & Y & 6/6 \\
meeting\_0\_story\_2 & Y & Y & Y & Y & Y & Y & 6/6 \\
meeting\_2\_story\_1 & Y & Y & Y & Y & Y & Y & 6/6 \\
meeting\_2\_story\_2 & Y & Y & Y & Y & P & Y & 5/6 \\
meeting\_30\_story\_1 & Y & Y & Y & Y & Y & Y & 6/6 \\
meeting\_117\_story\_1 & Y & Y & Y & Y & Y & Y & 6/6 \\
\midrule
\textbf{Average} & \multicolumn{6}{c}{} & \textbf{5.8/6} \\
\bottomrule
\end{tabular}
\end{table}

Note: Y = Yes (fully meets criterion), P = Partial (partially meets criterion)

The majority of generated user stories (83\%) achieved a perfect INVEST score of 6/6. Stories that scored 5/6 typically violated the ``Small'' criterion, indicating they could be further decomposed into more granular stories.

\subsubsection{Qualitative Evaluation Framework}

To systematically evaluate the generated outputs, we developed a qualitative assessment framework with the following criteria:

\begin{table}[H]
\centering
\caption{Qualitative evaluation criteria and results}
\label{tab:qualitative_eval}
\begin{tabular}{lcc}
\toprule
\textbf{Criterion} & \textbf{Score} & \textbf{Assessment} \\
\midrule
Persona Identification Accuracy & 10/10 & Excellent \\
Characteristic Quality & 8/10 & Very Good \\
User Story Format Compliance & 10/10 & Excellent \\
Contextualization & 9/10 & Excellent \\
Excerpt Selection Relevance & 8/10 & Very Good \\
Metadata Completeness & 10/10 & Excellent \\
Global Coherence & 8/10 & Very Good \\
\midrule
\textbf{Overall Average} & \textbf{9/10} & \textbf{Excellent} \\
\bottomrule
\end{tabular}
\end{table}

\subsubsection{Identified Strengths}

Key strengths include: 100\% accuracy in persona count detection, consistent user story format adherence, adequate contextualization with relevant characteristics, traceability through dialogue excerpts, and complete metadata in all outputs.

\subsubsection{Identified Limitations}

Limitations include: UTF-8 encoding issues with special characters, some overly broad user stories requiring decomposition, and occasional role overlap (e.g., ``Facilitator'' vs ``Moderator'').


\section{Conclusions and Future Work}
\label{sec:conclusion}

This work presented a multi-agent system for automated requirements elicitation using Generative AI. Two of three planned agents were successfully implemented: the Transcription Agent demonstrated reliable language detection and transcription quality, while the Identification Agent achieved excellent results in persona identification and INVEST-compliant user story generation.

\subsection{Limitations}

\begin{enumerate}
    \item Limited evaluation scope for Agent 1 (5 audio samples, no WER metrics)
    \item Incomplete system (Agent 3 not implemented)
    \item Small validation sample (7 dialogues)
    \item UTF-8 encoding issues requiring correction
    \item Some user stories violate INVEST "Small" criterion
    \item No real-world validation with software development teams
\end{enumerate}

\subsection{Future Work}

The second phase of this project (PFC2) will focus on completing the system implementation and conducting comprehensive quantitative evaluation. The planned activities include:

\subsubsection{Agent 3 Implementation}

Complete the development of the Formatting Agent, responsible for reading user stories generated by Agent 2 and producing standardized requirements documentation. The specific format and structure of the output document will be defined during PFC2, potentially including formats such as PDF, Markdown, or integration with requirements management tools.

\subsubsection{Agent 1 Enhancements and Evaluation}

Improve the Transcription Agent through the following activities:

\begin{itemize}
    \item Implement continuous audio recording capability, allowing users to record entire meetings without manual intervention
    \item Conduct rigorous evaluation using publicly available audio datasets with ground truth transcriptions (e.g., Common Voice, LibriSpeech)
    \item Calculate quantitative metrics including Word Error Rate (WER), Character Error Rate (CER), and processing time benchmarks
    \item Evaluate performance across different audio quality levels, background noise conditions, and multiple speakers scenarios
    \item Extend language support beyond English and Portuguese
\end{itemize}

\subsubsection{Agent 2 Enhancements and Evaluation}

Enhance the Identification Agent's robustness and expand evaluation scope:

\begin{itemize}
    \item Test persona identification with unstructured text without speaker labels, increasing identification difficulty and validating prompt engineering effectiveness
    \item Expand validation dataset to 50-100 meetings from AMI Corpus and other sources
    \item Resolve UTF-8 character encoding issues affecting special character representation
    \item Refine user story generation to improve granularity and ensure consistent compliance with INVEST criteria
    \item Implement automated detection of overly broad user stories requiring decomposition
\end{itemize}

\subsubsection{Comprehensive Quantitative Evaluation}

Replace the current qualitative assessment with rigorous quantitative metrics:

\begin{itemize}
    \item \textbf{Agent 1}: Word Error Rate (WER), Character Error Rate (CER), Real-Time Factor (RTF)
    \item \textbf{Agent 2}: Precision, Recall, and F1-score for persona identification; INVEST criteria compliance rate; semantic similarity metrics for user story quality
    \item \textbf{System-wide}: End-to-end accuracy from audio input to final requirements document; processing time per meeting; comparison with baseline manual methods
\end{itemize}

\subsubsection{End-to-End System Validation}

Conduct complete pipeline testing to evaluate the integrated system:

\begin{itemize}
    \item Execute full workflow from audio recording (Agent 1) through persona identification (Agent 2) to requirements document generation (Agent 3)
    \item Measure accuracy degradation across pipeline stages
    \item Identify error propagation patterns between agents
    \item Validate final output quality against manually created requirements documents
\end{itemize}

\subsubsection{Real-World Validation}

Evaluate the system in authentic software development contexts:

\begin{itemize}
    \item Conduct field tests with actual software development teams
    \item Compare system-generated requirements with manually elicited requirements
    \item Gather user feedback on usability and practical utility
    \item Assess time savings and cost reduction compared to traditional methods
    \item Identify real-world challenges not present in controlled dataset evaluation
\end{itemize}

These future activities will provide comprehensive validation of the proposed multi-agent system and enable assessment of its viability for practical deployment in software development organizations.


\bibliographystyle{sbc}
\bibliography{references}

\end{document}